\documentclass[11pt]{article}
\usepackage{amsmath, amssymb}
\usepackage[margin=1in]{geometry}

\title{Random Sea State Profile Generation (Deep Water)}
\author{}
\date{}

\begin{document}
\maketitle

\section{Domain}

Periodic domain of length $L = 164$, discretized on $N_x = 1024$ grid points. Gravity $g = 1$ (non-dimensional). Depth $h = 1000$ (effectively infinite).

\section{Surface elevation $\eta$}

Each sample is a band-limited Gaussian random field in Fourier space. For each sample, three hyperparameters are drawn uniformly:
\begin{itemize}
    \item Significant wave height: $H_s \sim \mathcal{U}(0.1, 0.6)$
    \item Peak wavenumber: $k_p \sim \mathcal{U}(0.06, 0.50)$
    \item Bandwidth parameter: $b_w \sim \mathcal{U}(1.5, 5.0)$
\end{itemize}

The spectral amplitude at wavenumber $k$ is a Gaussian centered on the peak wavenumber:
\begin{equation}
    S(k) = \exp\!\left(-\frac{(|k| - k_p)^2}{2\sigma_k^2}\right), \qquad \sigma_k = \frac{k_p}{b_w}.
\end{equation}

Fourier coefficients are constructed as $\hat{\eta}_k = \sqrt{S(k)}\, e^{i\phi_k}$ with uniformly random phases $\phi_k \sim \mathcal{U}(0, 2\pi)$, enforcing conjugate symmetry $\hat{\eta}_{-k} = \overline{\hat{\eta}_k}$ for a real-valued field. The result is then rescaled so that
\begin{equation}
    \operatorname{std}(\eta) = \frac{H_s}{4},
\end{equation}
matching the standard relation between significant wave height and surface standard deviation.

\section{Surface potential $\xi$}

The surface potential is derived from $\eta$ via the linear deep-water dispersion relation. In Fourier space:
\begin{equation}
    \hat{\xi}_k = -i\, \operatorname{sign}(k) \sqrt{\frac{g}{|k|}}\; \hat{\eta}_k.
\end{equation}
This is the exact relationship for linear deep-water waves where $\omega(k) = \sqrt{g|k|}$, giving the velocity potential at the surface consistent with the surface elevation. The zero mode is excluded ($\hat{\xi}_0 = 0$).

\section{DNO target $G(\eta;\, h)\,\xi$}

The Dirichlet--Neumann operator is evaluated using a truncated Taylor series in $\eta$ (order $M = 6$) via \texttt{dno\_series\_eval}, with a dealiasing pad factor of 8. This gives the vertical fluid velocity at the free surface corresponding to the $(\eta, \xi)$ pair at depth $h = 1000$, exact to truncation order.

Evaluation is batched (batch size 64) to exploit JAX's vectorized FFT operations.

\section{Dataset statistics (100k samples)}

\begin{center}
\begin{tabular}{ll}
    $\|\eta\|_\infty$ & 0.68 \\
    $\|\xi\|_\infty$ & 1.52 \\
    $\|G(\eta)\xi\|_\infty$ & 0.46
\end{tabular}
\end{center}

\section{Motivation}

Deep water does not support large-amplitude localized solitary gravity waves (a classical result). Stokes waves are the canonical nonlinear deep-water profiles but have limited spectral diversity (single-frequency with harmonics). Random sea states provide a physically realistic, spectrally diverse training distribution that covers the smooth, moderate-amplitude regime relevant to deep-water wave dynamics.

\end{document}
